% BOOK-FILE-ID: f-6dfda5da3083
% subsubsection: 極形式による図形的解釈

\subsection{複素平面で見る}\label{節：複素平面で見る}
  面白い話をしよう.
  一見複雑な式変形が,視覚的に分かりやすいような操作として解釈できるというのは,大きな驚きと喜びを持つ.
  ここでは\sectionref*{節：一次分数型}の複雑な式変形を図形的に解釈する.
  複素平面の上で漸化式を眺めてみよう.\footnote{複素平面や極形式に自信がない場合は\sectionref*{節：複素平面と極形式}を参照しよう.}

  結論から述べれば,この変形によって漸化式を「回転」と「拡大・縮小」として解釈することができる.
 
  $p,q,r,s$が実数であるとき,
  \[
    rx^2+(s-p)x-q=0
  \]
  の2解が実数でなければ,2解は互いに\bookterm{conjugate}{共役}な複素数となる.
  そこで,その2解を
  \[
    \beta=\overline{\alpha}
  \]
  とする.
  
  ここでは$b_n$と後に用いる$\lambda$を複素数とする.このとき,
  \[
    b_n=\frac{a_n-\alpha}{a_n-\overline{\alpha}}
  \]
  と置くと,\sectionref*{節：一次分数型}と同じ変形によって
  \[
    b_{n+1}
    =
    \frac{p-r\alpha}{p-r\overline{\alpha}}b_n
  \]
  となる.
  
  ここで,
  \[
    \lambda=\frac{p-r\alpha}{p-r\overline{\alpha}}
  \]
  とおく.
  $p,r$は実数であり,$\overline{p-r\alpha}=p-r\overline{\alpha}$である.
  分子と分母は互いに共役で絶対値が等しいから,
  \[
    |\lambda|
    =
    \left|
    \frac{p-r\alpha}{p-r\overline{\alpha}}
    \right|
    =1
  \]
  となる.
  
  したがって,$\lambda$は極形式を用いて
  \[
    \lambda=\cos\theta+i\sin\theta
  \]
  と表すことができる.
  
  一方,$b_n$も極形式を用いて
  \[
    b_n=\rho_n(\cos\varphi_n+i\sin\varphi_n)
  \]
  と表すことができる.
  このとき,
  \begin{align*}
    b_{n+1}
    &=
    (\cos\theta+i\sin\theta)
    \rho_n(\cos\varphi_n+i\sin\varphi_n)\\
    &=
    \rho_n
    \{\cos(\varphi_n+\theta)
    +i\sin(\varphi_n+\theta)\}
  \end{align*}
  であるから,
  \[
    \rho_{n+1}=\rho_n,
    \qquad
    \varphi_{n+1}=\varphi_n+\theta
  \]
  となる.
  
  つまり,$b_n$は複素平面上で原点からの距離を変えずに,$\theta$だけずつ回転している.
  
  このことから,
  \[
    b_{n+1}=\lambda b_n
  \]
  という等比数列への帰着は,単なる計算上の技巧ではなく,複素平面上の図形的な変化を表していることが分かる.
  
  実際に,次の漸化式を考えてみよう.
  \[
    a_{n+1}=\frac{a_n+1}{1-a_n}
  \]
  実数の初項が$1,0,-1$のときは途中で分母が0になり,数列を続けられない.
  以下の例では$a_1=\frac12$とするので,この問題は起こらない.
  この漸化式の不動点は
  \[
    x=\frac{x+1}{1-x}
  \]
  を解けばよいから,
  \[
    x^2+1=0
  \]
  より
  \[
    \alpha=i,\qquad\beta=-i
  \]
  である.
  
  そこで,
  \[
    b_n=\frac{a_n-i}{a_n+i}
  \]
  と置くと,
  \begin{align*}
    b_{n+1}
    &=
    \frac{1+i}{1-i}b_n
  \end{align*}
  となる.
  
  ここで,
  \[
    \frac{1+i}{1-i}
    =
    \frac{(1+i)^2}{2}
    =i
  \]
  であるから,
  \[
    b_{n+1}=ib_n
  \]
  を得る.
  
  $i$は極形式で
  \[
    i=\cos\frac{\pi}{2}+i\sin\frac{\pi}{2}
  \]
  と表せる.
  したがって,$b_n$は複素平面上で毎回反時計回りに$\frac{\pi}{2}$だけ回転する.
  
  例えば,$a_1=\frac12$とすると,
  \begin{align*}
    a_1&=\frac12\\
    a_2&=3\\
    a_3&=-2\\
    a_4&=-\frac13\\
    a_5&=\frac12
  \end{align*}
  となり,$a_n$は
  \[
    \frac12\longrightarrow3\longrightarrow-2
    \longrightarrow-\frac13\longrightarrow\frac12\longrightarrow\cdots
  \]
  と変化する.
  
  一方,$b_n$に着目すると,
  \[
    b_{n+1}=ib_n
  \]
  であるから,
  \[
    b_1\longrightarrow ib_1\longrightarrow-b_1
    \longrightarrow-ib_1\longrightarrow b_1\longrightarrow\cdots
  \]
  となる.
  
  つまり,実数上では複雑に見える
  \[
    a_1,a_2,a_3,\ldots
  \]
  の変化が,
  \[
    b_n=\frac{a_n-i}{a_n+i}
  \]
  という変換によって,複素平面上で一定角度ずつ回転する点の運動として表される.
  
  ここで,不動点$\alpha,\beta$が異なる実数の場合も見てみよう.
  係数$p,q,r,s$が実数なら,
  \[
    b_n=\frac{a_n-\alpha}{a_n-\beta},
    \qquad
    \lambda=\frac{p-r\alpha}{p-r\beta}
  \]
  とおいたとき,公比$\lambda$も実数であり,同じように
  \[
    b_{n+1}=\lambda b_n
  \]
  となる.
  $\lambda>0$なら原点からの距離を$\lambda$倍する.
  $\lambda<0$なら,距離を$|\lambda|$倍すると同時に,原点を中心に$\pi$だけ回転する.
  特に,$\lambda=-1$なら拡大・縮小を伴わず,半回転だけになる.

  実際に,
  \[
    a_1=2,\qquad a_{n+1}=\frac{1}{a_n}
  \]
  では,不動点は$\alpha=1,\ \beta=-1$である.
  $a_n$は$2$と$\frac{1}{2}$を交互に取るので,元の漸化式も次の変換も分母が0になることはない.
  \[
    b_n=\frac{a_n-1}{a_n+1}
  \]
  と置けば,公比は$\lambda=-1$となり,
  \[
    b_{n+1}=-b_n
  \]
  を得る.これは,実数の不動点を持つ場合でも,伸縮を伴わない回転として解釈できる例である.

  このように,差の比への変換は,公比の偏角に応じた回転と,公比の絶対値に応じた拡大・縮小を表している.
  係数が実数で不動点が実数でない場合の特徴は,公比の絶対値が自動的に1となり,拡大・縮小を伴わない回転になることである.

  \begin{bookpoint}
    \textbf{POINT}\par
    差の比を等比型へ移すと,公比の偏角が回転角,絶対値が距離の倍率になる.
    実数の負の公比でも半回転を伴い,実数係数で不動点が共役な虚数なら公比の絶対値は1となる.
  \end{bookpoint}
